\documentclass[12pt, twoside]{article}
%\documentclass[12pt, twoside]{article}
\usepackage[letterpaper, margin=1in, headsep=0.2in]{geometry}
\setlength{\headheight}{0.6in}
%\usepackage[english]{babel}
\usepackage[utf8]{inputenc}
\usepackage{microtype}
\usepackage{amsmath}
\usepackage{amssymb}
%\usepackage{amsfonts}
\usepackage[nomessages]{fp} %\FPeval{\var-name}{2*sin(pi/6)}
\usepackage{siunitx} %units in math. eg 20\milli\meter
\usepackage{yhmath} % for arcs, overparenth command
\usepackage{tikz} %graphics
\usetikzlibrary{quotes, angles, arrows, arrows.meta}
\usepackage{pgfplots}
\usepgfplotslibrary{fillbetween}
\pgfplotsset{compat=1.16}
\usepackage{graphicx} %consider setting \graphicspath{{images/}}
\usepackage{parskip} %no paragraph indent
\usepackage{enumitem}
\usepackage{multicol}
\usepackage{venndiagram}

\usepackage{fancyhdr}
\pagestyle{fancy}
\fancyhf{}
\renewcommand{\headrulewidth}{0pt} % disable the underline of the header
\raggedbottom
\hfuzz=2mm %suppresses overfull box warnings

\usepackage{hyperref}

\fancyhead[LO]{La Scuola d'Italia / Dr.\ Huson / IB Math 12 \\* 6 October 2026}
\fancyhead[RO]{\fbox{\begin{minipage}{6cm}\footnotesize Nickname:\\[0.5cm]\end{minipage}}}
\fancyhead[LE]{\fbox{\begin{minipage}{6cm}\footnotesize Nickname:\\[0.5cm]\end{minipage}}}
\fancyfoot[C]{\thepage}

\begin{document}
\subsubsection*{1.8 Classwork: Standard deviation on the GDC}

\emph{Use your GDC. Give answers exactly, or correct to three significant
figures. Write down the values you put into the calculator.}

\begin{enumerate}[itemsep=0.15cm]

%--- Problem 1: Mean, sd and variance of raw data; effect of a linear transformation (C to F) on mean and sd [9 marks] ---
%--- Source: Original (freshly authored, AA SL 4.3), GDC One-Variable Statistics from a list ---
\item[1.]

The maximum daily temperature, in degrees Celsius, was recorded in Rome on
ten days in October.

\begin{center}
\setlength{\tabcolsep}{8pt}\renewcommand{\arraystretch}{1.2}
\begin{tabular}{|l|c|c|c|c|c|c|c|c|c|c|}
\hline
Temperature ($^{\circ}$C) & 18 & 21 & 19 & 23 & 22 & 20 & 17 & 24 & 21 & 25 \\
\hline
\end{tabular}
\end{center}

\begin{enumerate}[itemsep=0.1cm]
  \item Find the mean and the standard deviation of these temperatures.
  \item Write down the variance of these temperatures.
\end{enumerate}

A visitor from New York converts each temperature to degrees Fahrenheit
using the formula $F = 1.8\,C + 32$.

\begin{enumerate}[resume, itemsep=0.1cm]
  \item Find the mean of the ten temperatures in $^{\circ}$F.
  \item Find the standard deviation of the ten temperatures in $^{\circ}$F.
  \item Explain why the ``$+\,32$'' in the formula changes the mean but
  does not change the standard deviation.
\end{enumerate}

\begin{tikzpicture}
  \draw (-0.6,0) rectangle (14.6,11.5);
  \foreach \y in {1,2,...,11} \draw[dotted] (0,\y)--(14.1,\y);
\end{tikzpicture}

\newpage

%--- Problem 2: Frequency table with an unknown frequency from a given mean; sd; values beyond one sd; effect of an added value [7 marks] ---
%--- Source: Original (freshly authored, AA SL 4.3), GDC One-Variable Statistics with a frequency list ---
\item[2.]

The table shows the number of goals, $x$, that a football team scored in
each of its matches last season. The mean number of goals per match was
$2.2$.

\begin{center}
\setlength{\tabcolsep}{10pt}\renewcommand{\arraystretch}{1.2}
\begin{tabular}{|l|c|c|c|c|c|c|}
\hline
Number of goals, $x$ & 0 & 1 & 2 & 3 & 4 & 5 \\
\hline
Frequency, $f$ & 2 & 7 & $k$ & 6 & 4 & 1 \\
\hline
\end{tabular}
\end{center}

\begin{enumerate}[itemsep=0.1cm]
  \item Find the value of $k$.
  \item Find the standard deviation of the number of goals per match.
  \item Find the number of matches in which the number of goals was more
  than one standard deviation away from the mean.
  \item In its first match of the new season the team scores $2$ goals.
  This match is added to the data. State whether the mean and the
  standard deviation increase, decrease or stay the same. Give a reason
  for each answer.
\end{enumerate}

\begin{tikzpicture}
  \draw (-0.6,0) rectangle (14.6,14.5);
  \foreach \y in {1,2,...,14} \draw[dotted] (0,\y)--(14.1,\y);
\end{tikzpicture}

\end{enumerate}
\end{document}
